% Antipodal ear insertion: new vertices u (after j) and u' (after j+h); the new diameter is (u, u').
\begin{figure}[H]
\centering
\begin{tikzpicture}[scale=0.85, every node/.style={font=\footnotesize}]
  \foreach \k/\lab in {0/{V_j},1/{V_{j+1}},2/{},3/{},4/{},5/{V_{j+h}},6/{V_{j+h+1}},7/{},8/{},9/{}} {
    \coordinate (P\k) at ({90-36*\k}:2);
  }
  \draw[thick] (P0) -- (P1) -- (P2) -- (P3) -- (P4) -- (P5) -- (P6) -- (P7) -- (P8) -- (P9) -- cycle;
  \coordinate (U)  at ($(P0)!0.5!(P1) + ({72}:0.75)$);
  \coordinate (U2) at ($(P5)!0.5!(P6) + ({-108}:0.75)$);
  \draw[thick, orange!80!black] (P0) -- (U) -- (P1);
  \draw[thick, orange!80!black] (P5) -- (U2) -- (P6);
  \draw[dashed] (P0) -- (P1);
  \draw[dashed] (P5) -- (P6);
  \draw[very thick, red!70!black] (U) -- (U2);
  \draw[gray] (P0) -- (P5);
  \draw[gray] (P1) -- (P6);
  \foreach \k in {0,...,9} \fill (P\k) circle (1.5pt);
  \fill[orange!80!black] (U) circle (2pt);
  \fill[orange!80!black] (U2) circle (2pt);
  \node[above left] at (P0) {$V_j$};
  \node[right] at (P1) {$V_{j+1}$};
  \node[below right] at (P5) {$V_{j+h}$};
  \node[left] at (P6) {$V_{j+h+1}$};
  \node[above right] at (U) {$u = V_j + V_{j+1}$};
  \node[below left] at (U2) {$u' = V_{j+h} + V_{j+h+1}$};
\end{tikzpicture}
\caption{Antipodal ear insertion. The old diameters (grey: two of them) keep their values; the only new
diameter is $(u, u')$ (red), of value $S = \det(V_j + V_{j+1}, V_{j+h} + V_{j+h+1})$.}
\end{figure}
